% ============================================================================
% WORKBOOK — Section 8.7: Data Displays Extended (Additional)
% State-specific (AK, IN, OK)
% Practice-focused reinforcement for the study guide topic
% ============================================================================
\section{Data Displays Extended}
\topicTitle{Data Displays Extended}

\begin{quickReview}{Data Displays Extended}
Different displays suit different data questions:

\begin{itemize}[leftmargin=*]
    \item \textbf{Frequency table:} Lists categories and their counts. Good for exact values.
    \item \textbf{Double bar graph:} Compares two groups side by side using bars.
    \item \textbf{Line graph:} Shows how data changes over time. Points are connected.
    \item \textbf{Histogram:} Groups numerical data into intervals (bins). Shows the shape of a distribution.
\end{itemize}

\medskip
\textbf{Choosing the right display:}
Use a \textbf{line graph} for changes over time, a \textbf{bar graph} to compare categories, a \textbf{double bar graph} to compare two groups, and a \textbf{frequency table} for exact counts.
\end{quickReview}

% ── Warm-Up ──────────────────────────────────────────────────────────────────
\begin{practiceBox}[funGreen]{\faSun~Warm-Up}

\practiceHeader[funGreen]{Choose the Best Display}

\prob You want to show how the temperature changed each day for a week. Which display? \answerBlank[3cm]
\answerExplain{Line graph}{A line graph shows how data changes over time, making it ideal for daily temperature trends.}

\prob You want to compare the number of boys and girls in each grade. Which display? \answerBlank[3cm]
\answerExplain{Double bar graph}{A double bar graph compares two groups (boys and girls) across categories (grades) side by side.}

\prob You want to show how many students scored in each point range on a test ($60$--$69$, $70$--$79$, etc.). Which display? \answerBlank[3cm]
\answerExplain{Histogram}{A histogram groups numerical data into intervals and shows the distribution shape.}

\prob You want to list exactly how many students chose each lunch option. Which display? \answerBlank[3cm]
\answerExplain{Frequency table}{A frequency table gives exact counts for each category.}

\prob You want to track your savings account balance each month for a year. Which display? \answerBlank[3cm]
\answerExplain{Line graph}{Savings over time is a change-over-time scenario, which is best shown with a line graph.}

\end{practiceBox}

% ── Practice A: Reading Frequency Tables ─────────────────────────────────────
\begin{practiceBox}{Reading Frequency Tables}

\practiceHeader{Use the frequency table to answer the questions.}

\begin{center}
\begin{tabular}{lc}
    \textbf{Favorite Subject} & \textbf{Frequency} \\
    \hline
    Math     & $15$ \\
    Science  & $12$ \\
    English  & $9$ \\
    History  & $6$ \\
    Art      & $8$ \\
\end{tabular}
\end{center}

\prob How many students were surveyed in total? \answerBlank[3cm]
\answerExplain{$50$ students}{$15 + 12 + 9 + 6 + 8 = 50$. Add all the frequencies together.}

\prob What fraction of students chose Science? \answerBlank[3cm]
\answerExplain{$\frac{12}{50} = \frac{6}{25}$}{$12$ out of $50$ students chose Science. Simplified: $\frac{12}{50} = \frac{6}{25}$.}

\prob What percent of students chose Art? Round to the nearest whole percent. \answerBlank[3cm]
\answerExplain{$16\%$}{$\frac{8}{50} = 0.16 = 16\%$.}

\prob Which subject was chosen by exactly $\frac{3}{25}$ of the students? \answerBlank[3cm]
\answerExplain{History}{$\frac{3}{25} = \frac{6}{50}$. History has a frequency of $6$, so $\frac{6}{50} = \frac{3}{25}$.}

\prob How many more students chose Math than English? \answerBlank[3cm]
\answerExplain{$6$ students}{$15 - 9 = 6$ more students chose Math.}

\end{practiceBox}

% ── Practice B: Double Bar Graphs ────────────────────────────────────────────
\begin{practiceBox}[funTeal]{Double Bar Graphs}

\practiceHeader[funTeal]{The table shows test scores for two classes. Use it to answer the questions.}

\begin{center}
\begin{tabular}{lccccc}
    & \textbf{Test 1} & \textbf{Test 2} & \textbf{Test 3} & \textbf{Test 4} & \textbf{Test 5} \\
    \hline
    Class A & $72$ & $78$ & $80$ & $85$ & $88$ \\
    Class B & $68$ & $75$ & $82$ & $79$ & $90$ \\
\end{tabular}
\end{center}

\prob On which test did Class B score higher than Class A? \answerBlank[3cm]
\answerExplain{Tests $3$ and $5$}{Test 3: B($82$) $>$ A($80$). Test 5: B($90$) $>$ A($88$). On all other tests, Class A scored higher.}

\prob What is the mean score for Class A across all $5$ tests? \answerBlank[3cm]
\answerExplain{$80.6$}{$\frac{72+78+80+85+88}{5} = \frac{403}{5} = 80.6$.}

\prob What is the mean score for Class B across all $5$ tests? \answerBlank[3cm]
\answerExplain{$78.8$}{$\frac{68+75+82+79+90}{5} = \frac{394}{5} = 78.8$.}

\prob Which class improved more from Test 1 to Test 5? \answerBlank[3cm]
\answerExplain{Class B}{Class A: $88 - 72 = 16$ points. Class B: $90 - 68 = 22$ points. Class B improved by $22$ points vs.\ $16$.}

\prob On how many tests was the difference between the two classes more than $3$ points? \answerBlank[3cm]
\answerExplain{$2$ tests}{Test 1: $|72 - 68| = 4 > 3$. Test 2: $|78 - 75| = 3$, not more. Test 3: $|80 - 82| = 2$. Test 4: $|85 - 79| = 6 > 3$. Test 5: $|88 - 90| = 2$. Two tests (1 and 4) had differences greater than $3$.}

\end{practiceBox}

% ── Practice C: Line Graphs ──────────────────────────────────────────────────
\begin{practiceBox}[funOrange]{Line Graphs — Changes Over Time}

\practiceHeader[funOrange]{A store tracked weekly sales (in hundreds of dollars):}

\begin{center}
\begin{tabular}{lccccc}
    \textbf{Week} & $1$ & $2$ & $3$ & $4$ & $5$ \\
    \hline
    \textbf{Sales (\$100s)} & $30$ & $35$ & $28$ & $40$ & $42$ \\
\end{tabular}
\end{center}

\prob Between which two weeks did sales increase the most? \answerBlank[3cm]
\answerExplain{Week $3$ to Week $4$}{Week 3 to 4: $40 - 28 = 12$. This is the biggest increase. (Week 1 to 2: $5$; Week 4 to 5: $2$).}

\prob Between which two weeks did sales decrease? \answerBlank[3cm]
\answerExplain{Week $2$ to Week $3$}{Sales went from $35$ to $28$, a decrease of $7$. This is the only week-to-week decrease.}

\prob What is the overall change from Week $1$ to Week $5$? \answerBlank[3cm]
\answerExplain{Increase of $\$1{,}200$}{$42 - 30 = 12$ (in hundreds). $12 \times 100 = \$1{,}200$ increase.}

\prob What is the average weekly sales amount? \answerBlank[3cm]
\answerExplain{$\$3{,}500$}{$\frac{30+35+28+40+42}{5} = \frac{175}{5} = 35$ (in hundreds) $= \$3{,}500$.}

\end{practiceBox}

% ── True or False ────────────────────────────────────────────────────────────
\begin{practiceBox}[funPurple]{True or False?}

\trueOrFalse{A bar graph is the best choice for showing change over time.}
\answerExplain{False}{A line graph is better for showing change over time because the connected points emphasize trends. Bar graphs compare separate categories.}

\trueOrFalse{In a frequency table, the sum of all frequencies equals the total number of data values.}
\answerExplain{True}{Each frequency counts how many data values fall in that category. The total of all frequencies is the total data count.}

\trueOrFalse{A double bar graph can compare exactly two groups.}
\answerExplain{True}{A double bar graph uses two bars side by side for each category, comparing two groups.}

\end{practiceBox}

% ── Word Problems ────────────────────────────────────────────────────────────
\begin{practiceBox}[funBlue]{\faGlobe~Word Problems}

\wordProblem{A school recorded absences for two grades over four months:

\begin{center}
\begin{tabular}{lcccc}
    & \textbf{Sep} & \textbf{Oct} & \textbf{Nov} & \textbf{Dec} \\
    \hline
    Grade 6 & $12$ & $15$ & $20$ & $25$ \\
    Grade 7 & $10$ & $18$ & $22$ & $30$ \\
\end{tabular}
\end{center}

Which grade had a larger increase in absences from September to December?}{grade}
\answerExplain{Grade 7}{Grade 6 increase: $25 - 12 = 13$. Grade 7 increase: $30 - 10 = 20$. Grade 7 had a larger increase.}

\wordProblem{A frequency table shows: Red $= 24$, Blue $= 18$, Green $= 12$, Yellow $= 6$. What percent of the total is Red?}{percent}
\answerExplain{$40\%$}{Total: $24 + 18 + 12 + 6 = 60$. Red: $\frac{24}{60} = 0.40 = 40\%$.}

\end{practiceBox}

% ── Challenge ────────────────────────────────────────────────────────────────
\begin{challengeBox}
\prob A company tracks monthly revenue and expenses (in thousands):

\begin{center}
\begin{tabular}{lcccccc}
    & \textbf{Jan} & \textbf{Feb} & \textbf{Mar} & \textbf{Apr} & \textbf{May} & \textbf{Jun} \\
    \hline
    Revenue  & $40$ & $45$ & $50$ & $48$ & $55$ & $60$ \\
    Expenses & $35$ & $40$ & $42$ & $50$ & $45$ & $48$ \\
\end{tabular}
\end{center}

In which month(s) did expenses exceed revenue? What was the total profit (revenue $-$ expenses) over the six months? \answerBlank[3cm]
\answerExplain{April; total profit $= \$38{,}000$}{April: expenses $50$ $>$ revenue $48$, so expenses exceeded revenue. Monthly profits: Jan $5$, Feb $5$, Mar $8$, Apr $-2$, May $10$, Jun $12$. Total: $5+5+8+(-2)+10+12 = 38$ thousand $= \$38{,}000$.}
\end{challengeBox}

\encouragement{Amazing work with data displays! You can now read, create, and compare many types of graphs and tables!}
